\documentclass[10pt,a4paper]{amsart}
\usepackage[margin=19mm]{geometry}
\usepackage{amsmath,amssymb,mathtools}
\usepackage[scr=rsfs,cal=euler]{mathalfa}
\usepackage{xfrac}
\usepackage[unicode,hidelinks,bookmarks=false]{hyperref}
\newcommand{\ie}{i.e.\ }
\newcommand{\cf}{cf.\ }
\newcommand{\resp}{respectively\ }
\newcommand{\Subsection}{\subsection}
\providecommand{\nobreakdash}{\nobreak\hbox{-}}
\setlength{\emergencystretch}{2em}
\multlinegap0pt
\usepackage{enumerate}
\usepackage{amssymb}
\usepackage{algorithm}\usepackage{algpseudocode}
\usepackage{tcolorbox}
\usepackage{float}
\usepackage{csquotes}
\newtheorem{thm}{Theorem}
\newtheorem{pf}{Pf}
\title{Piecewise Linear Approximation and PID Control Optimization for Nonlinear Systems}
\author{Robert Vrabel}
\date{Source: arXiv:2501.00820v2 (CC BY 4.0)}
\begin{document}
\maketitle

\noindent\emph{Source and licence.} Piecewise Linear Approximation and PID Control Optimization for Nonlinear Systems. Version arXiv:2501.00820v2 (CC BY 4.0), distributed by arXiv under the Creative Commons Attribution 4.0 International licence, http://creativecommons.org/licenses/by/4.0/, as recorded in the arXiv licence metadata for this exact version and shown on the abstract page https://arxiv.org/abs/2501.00820v2. Full licence notice: \texttt{licenses/arxiv-2501.00820-CC-BY-4.0.txt}.\par\smallskip

\noindent\textit{Editorial scope.} Counted unit: the article's thm thm (source0.tex lines 342--354). The article's complete original proof of it is delivered with it (source0.tex lines 355--488, one \texttt{proof} environment of 134 lines). No mathematics is written, repaired or re-derived: the statement and the proof are the article's own lines, in the article's own notation, and no retained line is substituted or re-flowed.  Retained uncounted as the source-local context the proof needs: lines 263--275.  Nothing was omitted from any retained span, so the package carries no omission record.  No retained line contains a citation, so the wrapper carries no bibliography and no bibliography entry.  No retained line contains a figure environment, an image-inclusion command or drawing code, so no image is delivered and the package declares no asset dependency.  Editorial formatting only, in addition: an AMS \texttt{amsart} preamble that restates the article's own declarations and macros used by the retained lines, namely the environments \texttt{thm}, \texttt{pf} on separate counters, together with the standard packages those lines need.  The retained environments print in this document as Theorem 1, Pf 1 in order of appearance, whereas the article prints the same items with the numbering of its own full document.  The complete native source of the article as distributed in the arXiv e-print for this exact version is bundled unchanged as \texttt{sources/171020/source0.tex}.
\begin{equation}\label{nonlinear_system_1_order}
\begin{aligned}
\left.
\begin{array}{ll}
y_1' &= y_2, \\
y_2' &= y_3, \\
&\ \vdots \\
y_{n-1}' &= y_n, \\
y_n' &= -f(y_1, y_2, \dots, y_n) + u(t)
\end{array}
\right\} \triangleq y'=\tilde f(y)+\tilde u(t).
\end{aligned}  
\end{equation}

\begin{thm} Let $f\in C^2(\mathbb{R}^n\mapsto \mathbb{R})$ is globally Lipschitz in \((y, y', \dots, y^{(n-1)})\) with Lipschitz constant $L$ and $u\in C(\mathbb{R}\mapsto \mathbb{R})$. Then the initial value problem
\[
y^{(n)} + f(y, y', \dots, y^{(n-1)}) = u(t),
\]
with a zero initial conditions
\[
y(0) = y'(0) = \dots = y^{(n-1)}(0) = 0,
\]
has a unique solution $y(t)$ defined on $\mathbb{R}$ satisfying here the inequality
\[
|y(t)|\leq Ke^{2\tilde L|t|},\quad  K\geq0,\quad \tilde L=({n-1+L^2})^{1/2}.
\]
\end{thm}

\begin{pf}
The \(n\)-th order differential equation can be reformulated as a system of first-order ODEs (\ref{nonlinear_system_1_order}).
On the space \( C(\mathbb{R} \to \mathbb{R}^n) \) of continuous functions, define the weighted supremum norm
\[
\|y\|_{\tilde L} = \sup_{t \in \mathbb{R}} e^{-2\tilde L|t|}\,|y(t)|.
\]
The subspace of functions with finite norm under this definition includes all constant functions. Including all constant functions in this space ensures that simple initial guesses, such as \(y_0(t) = 0\), lie within the domain of the Picard operator, allowing the iterative process to start consistently within the function space.
The Picard operator
\( (P(y))(t) \) with zero initial conditions is defined as
   \[
   (P(y))(t) = \int_0^t \big(\tilde f(y(\tau))+\tilde u(\tau) \big)\, d\tau,
   \]
   where \( \tilde f \) is a Lipschitz function with Lipschitz constant \( \tilde L > 0 \). This means
   \[
   |\tilde f(x) - \tilde f(y)| \leq \tilde L |x - y|, \quad \forall x, y\in \mathbb{R}^n.
   \]
To prove that \( (P(y))(t)) \) is contractive, we show that the mapping \( P: y(t) \mapsto (P(y))(t) \) satisfies the contraction condition, that is, there exists \( 0 \leq k < 1 \) such that
\[
\|P(x) - P(y)\|_{\tilde L} \leq k \|x - y\|_{\tilde L}.
\]
By definition
\[
(P(x))(t) - (P(y))(t) = \int_0^t \big(\tilde f(x(\tau)) - \tilde f(y(\tau))\big) \, d\tau.
\]
The norm is given by  
\[
\|P(x) - P(y)\|_{\tilde L} 
= \sup_{t \in \mathbb{R}} e^{-2\tilde L |t|} \left| \int_0^t \big(\tilde f(x(\tau)) - \tilde f(y(\tau))\big) \, d\tau \right|.
\]
By the Lipschitz condition on \( \tilde f \),  
\[
|\tilde f(x(\tau)) - \tilde f(y(\tau))| \leq \tilde L |x(\tau) - y(\tau)|.
\]  
This implies 
\[
\left| \int_0^t \big(\tilde f(x(\tau)) - \tilde f(y(\tau))\big) \, d\tau \right| \leq \int_0^t \tilde L |x(\tau) - y(\tau)| \, d\tau.
\]
Using the definition of the norm \( \|\cdot\|_{\tilde L} \),  
\[
|x(\tau) - y(\tau)| \leq e^{2\tilde L |\tau|} \|x - y\|_{\tilde L}.
\]  
Substituting this bound into the integral yields
\[
\left| \int_0^t \big(\tilde f(x(\tau)) - \tilde f(y(\tau))\big) \, d\tau \right| \leq \int_0^t \tilde L e^{2\tilde L |\tau|} \|x - y\|_{\tilde L} \, d\tau.
\]
To express the estimate in terms of the weighted supremum norm, multiply both sides by the factor \( e^{-2\tilde L |t|} \) 
\[
e^{-2\tilde L |t|} \left| \int_0^t \big(\tilde f(x(\tau)) - \tilde f(y(\tau))\big) \, d\tau \right| 
\leq e^{-2\tilde L |t|} \int_0^t \tilde L e^{2\tilde L |\tau|} \|x - y\|_{\tilde L} \, d\tau.
\]
For \( \tau \in [0, t] \), we have \( |\tau| = \tau \). Hence,  
\[
e^{-2\tilde L |t|} e^{2\tilde L |\tau|} = e^{-2\tilde L (t - \tau)}.
\]
Substituting this expression, we obtain 
\[
e^{-2\tilde L |t|} \left| \int_0^t \big(\tilde f(x(\tau)) - \tilde f(y(\tau))\big) \, d\tau \right| 
\leq \int_0^t \tilde L e^{-2\tilde L (t - \tau)} \|x - y\|_{\tilde L} \, d\tau.
\]
The exponential term \( e^{-2\tilde L (t - \tau)} \) integrates to a finite value 
\[
\int_0^t e^{-2\tilde L (t - \tau)} \, d\tau = \frac{1 - e^{-2\tilde L t}}{2\tilde L} \leq \frac{1}{2\tilde L}.
\]
Therefore,
\[
e^{-2\tilde L |t|} \left| \int_0^t \big(\tilde f(x(\tau)) - \tilde f(y(\tau))\big) \, d\tau \right| 
\leq  \frac{1}{2} \|x - y\|_{\tilde L}.
\]
For \( t < 0 \), the argument proceeds analogously, with only minor adjustments due to the sign of \( t \). These involve appropriately handling the integration limits and the exponential weighting term \( e^{-2\tilde L |t|} \), but the underlying reasoning remains identical.

Taking the supremum over \( t \in \mathbb{R} \), we obtain
\[
\|P(x)- P(y)\|_{\tilde L} \leq \frac{1}{2} \|x - y\|_{\tilde L}.
\]
Since \( \frac{1}{2} < 1 \), the operator \( P \) is a contraction. 

To apply the Banach Fixed-Point Theorem, we first need to verify that the Picard operator \( P(y) \) maps the chosen function space into itself. Specifically, we must show that if \( y(t) \) belongs to the space of continuous functions endowed with the weighted supremum norm \( \|y\|_{\tilde L} = \sup_{t \in \mathbb{R}} e^{-2\tilde L |t|} |y(t)| < \infty \), then \( P(y)(t) \) also lies in the same space.

Define the space \( \mathcal{C}_{\tilde L} \) as
   \[
   \mathcal{C}_{\tilde L} = \{ y \in C(\mathbb{R} \mapsto \mathbb{R}^n) \; : \; \|y\|_{\tilde L} < \infty \}.
   \]
A function \( y(t) \in \mathcal{C}_{\tilde{L}} \) satisfies
\[
\|y\|_{\tilde{L}} = \sup_{t \in \mathbb{R}} e^{-2\tilde{L} |t|} |y(t)| < \infty.
\]
This implies that, for all \( t \in \mathbb{R} \),
\[
|y(t)| \leq K e^{2\tilde{L}|t|},
\]
where \( K = \|y\|_{\tilde{L}} \) is a finite constant.

Recall that the Picard operator is defined by
\[
(P(y))(t) = \int_0^t \big(\tilde{f}(y(\tau)) + \tilde u(\tau)\big) \, d\tau,
\]
where \( \tilde{f} \) satisfies the Lipschitz condition with constant \( \tilde{L} > 0 \), and \( \tilde u(t) \in \mathcal{C}_{\tilde L} \) is a continuous input function.  
We aim to show that
\[
\text{if } y \in \mathcal{C}_{\tilde L}, \text{ then } P(y) \in \mathcal{C}_{\tilde L}.
\]

Since \( y \in \mathcal{C}_{\tilde L} \), it follows that \( \|\tilde{f}(y)\|_{\tilde L} < \infty \) because \( \tilde{f} \) is Lipschitz continuous, and thus
\[
\|\tilde{f}(y)\|_{\tilde L} \leq |\tilde{f}(0)| + \tilde{L} \|y\|_{\tilde L}.
\]
Moreover, \( \tilde u(t) \in \mathcal{C}_{\tilde L} \) is continuous and satisfies \( \|\tilde u\|_{\tilde L} < \infty \).  
Hence, the integrand \( \tilde{f}(y(\tau)) + \tilde u(\tau) \) is well defined and integrable for all \( t \in \mathbb{R} \).
To complete the proof, we need to verify that
\[
\|P(y)\|_{\tilde L} = \sup_{t \in \mathbb{R}} e^{-2\tilde L |t|} 
\left|\int_0^t \big(\tilde{f}(y(\tau)) + \tilde u(\tau)\big) \, d\tau \right| < \infty.
\]
Using the triangle inequality, we obtain
\[
\left|\int_0^t \big(\tilde{f}(y(\tau)) + \tilde u(\tau)\big) \, d\tau \right| 
\leq \int_0^t |\tilde{f}(y(\tau))| \, d\tau + \int_0^t |\tilde u(\tau)| \, d\tau.
\]
Applying the weighted supremum norm to both sides yields
\[
e^{-2\tilde{L}|t|} \left|\int_0^t \big(\tilde{f}(y(\tau)) + \tilde u(\tau)\big) \, d\tau \right| 
\leq e^{-2\tilde{L}|t|} \int_0^t e^{2\tilde{L}|\tau|} \|\tilde{f}(y)\|_{\tilde L} \, d\tau 
+ e^{-2\tilde{L}|t|} \int_0^t e^{2\tilde{L}|\tau|} \|\tilde u\|_{\tilde L} \, d\tau.
\]
Consequently,
\[
e^{-2\tilde{L}|t|} \left|\int_0^t \big(\tilde{f}(y(\tau)) + \tilde u(\tau)\big) \, d\tau \right|
\leq \|\tilde{f}(y)\|_{\tilde L} \cdot \frac{1}{2\tilde{L}} (1 - e^{-2\tilde{L}|t|}) +\|\tilde u\|_{\tilde L} \cdot \frac{1}{2\tilde{L}} (1 - e^{-2\tilde{L}|t|}).
\]
Taking the supremum over \( t \in \mathbb{R} \), we conclude that the norm is finite,
\[
\|P(y)\|_{\tilde L} \leq \frac{\|\tilde{f}(y)\|_{\tilde L} +\|\tilde u\|_{\tilde L} }{2\tilde{L}} < \infty.
\]
\end{pf} 

\end{document}
